\documentclass[11pt]{article}
\usepackage[margin=1in]{geometry}
\usepackage{amsmath,amssymb,amsthm,mathtools}
\usepackage[T1]{fontenc}
\usepackage{lmodern,microtype}
\usepackage[hidelinks]{hyperref}
\newtheorem{theorem}{Theorem}
\newcommand{\C}{\mathbb C}
\DeclareMathOperator{\diag}{diag}
\title{Which fixed conics suffice to test the order-four\\ Marcus--de Oliveira conjecture?}
\date{}
\author{}
\begin{document}
\maketitle
For a fixed relative unitary $W$, the difference form of the conjecture asks whether
\[
 \det(D_a-WD_cW^*)\ \in\
 \operatorname{conv}_{\mathbb R}\left\{\prod_i(a_i-c_{\pi(i)}):\pi\in S_n\right\},
 \qquad D_a=\diag(a_1,\ldots,a_n).
\]
Here $c=-b$ for the usual sum convention.
We show that testing only first spectra on the ellipse of semiaxes
$2+\sqrt3$ and $1$ already tests every order-four pair. We also give
a reduction to any one fixed nondegenerate parabola or hyperbola in
orders four and five. These statements concern sufficiency of the test
families. Positivity on the test families remains to be proved.

\section*{Projective transport}
If $f(z)=(Az+B)/(Cz+D)$ has no spectral pole, functional calculus gives
\[
 (CX+DI)(f(X)-f(Y))(CY+DI)=(AD-BC)(X-Y).
\]
Thus the determinant and all comparison products acquire the same nonzero
complex factor. The relative eigenbasis stays fixed. Given four distinct
points, any other ordered quadruple of the same cross-ratio is obtained
by a fractional linear map: mapping the first three points specifies the
map uniquely and the cross-ratio recovers the fourth. A required map can
have a pole at a partner eigenvalue. Approximate the partner by lists
avoiding the pole, apply the result in the original chart and pass to
the limit. Convex membership is closed by compactness of the simplex
of permutation weights.

\section*{The ellipse threshold}
For $r\ge1$ let $E_r=\{r\cos t+i\sin t:t\in\mathbb R\}$.
For distinct ordered points put
\[
 \lambda=\frac{(z_1-z_3)(z_2-z_4)}{(z_1-z_4)(z_2-z_3)}.
\]
Let $\Theta$ be the maximum absolute principal argument of the values of
$\lambda$ under relabeling, using $\pi$ for a negative real value.
\begin{theorem}
A quadruple is projectively equivalent, up to ordering, to four points of
$E_r$ precisely when $\Theta\ge4\arctan(1/r)$. Every quadruple is so
represented exactly when $r\ge 2+\sqrt3$.
\end{theorem}
\begin{proof}
First construct the points. Write $T(x+iy)=rx+iy$ and $e_s=e^{is}$.
Take $u=T(e_{s+d}), v=T(e_{s-d})$ with $0<d<\pi/2$. The
cross-ratio of $(u,-u,v,-v)$ is the square of
\[
 \frac{u-v}{u+v}=\tan d\,\frac{T(e_{s+\pi/2})}{T(e_s)} .
\]
The oriented angle $\theta$ of the latter quotient has
$\cot\theta=-(r^2-1)\sin(2s)/(2r)$ by the dot product and the
positive area $r$ of the two vectors. Its range is
$[2\arctan(1/r),\pi-2\arctan(1/r)]$. A cross-ratio with absolute
argument at least $4\arctan(1/r)$ has a square root with absolute
argument in $[2\arctan(1/r),\pi/2]$. Choose $s$ to give this
angle and $d$ to give the modulus. Conjugating the four points gives
the negative angle if needed.

For necessity the circle $r=1$ has only real cross-ratios, all of
which admit a negative ordering. For $r>1$, set $q=(r-1)/(r+1)$.
Up to a uniform scale every ellipse point has the form
$z_j=w_j+q/w_j$, $|w_j|=1$. Reorder so that the cross-ratio
$\lambda(w)$ on the circle is negative. Let $p=\prod_jw_j$,
choose either square root of $p$, and put $Z=\sqrt p/q+q/\sqrt p$.
Factoring $z_i-z_j=(w_i-w_j)(1-q/(w_iw_j))$ gives
\[
 \frac{\lambda(z)}{\lambda(w)}=\frac{Z-s}{Z-t},\qquad
 s=\frac{w_1w_3+w_2w_4}{\sqrt p},\quad
 t=\frac{w_1w_4+w_2w_3}{\sqrt p}.
\]
The paired summands are conjugates, so $s,t\in[-2,2]$.
The point $Z$ runs on the confocal ellipse with foci $\pm2$ and
semiaxes $A=q^{-1}+q$, $B=q^{-1}-q$. The angle subtended by
$[s,t]$ is bounded by that of $[-2,2]$. If the focal distances are
$h,k$, then $h+k=2A$ and the cosine of this last angle is
$2B^2/(hk)-1$. The angle is maximal at $h=k$, giving
$2\arctan(2/B)=4\arctan q$. Hence $|\arg\lambda(z)|\ge
\pi-4\arctan q=4\arctan(1/r)$ for the chosen ordering.

For nonreal $\lambda$, the arguments in absolute value of $\lambda$,
$1-\lambda$, and $(\lambda-1)/\lambda$ are the three interior angles of
the triangle $0,1,\lambda$. All three belong to its relabeling orbit.
The largest is at least $\pi/3$. The equilateral shape
$\lambda=e^{i\pi/3}$ has only $e^{\pm i\pi/3}$ in its orbit.
Real shapes have $\Theta=\pi$. Since $\cot(\pi/12)=2+\sqrt3$,
these facts give the asserted sharp threshold.
\end{proof}

If MOC holds on one ellipse at or above this threshold with arbitrary
partner and fixed $W$, transport proves MOC for that $W$ in order four.
Repeated points follow by closure. The threshold is sharp for
representation of quadruples. It gives no obstruction to MOC on a
thinner ellipse.

\section*{Open conics in orders four and five}
\begin{theorem}
In either order, testing every distinct first spectrum on one fixed
real nondegenerate parabola, or on one fixed real nondegenerate
hyperbola, suffices for MOC with a fixed $W$ and arbitrary partner.
\end{theorem}
\begin{proof}
For a generic five-point set choose three points on a circle whose
open disk omits the other two. Indeed a lower triangular face of the
convex hull of the lifts $(x,y,x^2+y^2)$ gives such a circle by
completing the square in its supporting-plane inequality. Generic
points have no collinear triple and no four concyclic. For four
points add a fifth generic auxiliary point and discard it after the
construction. A fractional linear map takes the chosen circle to the
real line and its exterior to the upper half-plane, with the three
marked boundary points finite. Write $u,v$ for the two upper points.

Apply a real fractional linear map $h_t(z)=-1/(z-t)$ with $t$ real.
The boundary triple remains real. The direction of the line through
$h_t(u),h_t(v)$ is that of $(u-v)/[(u-t)(v-t)]$. As $t$ increases
from $-\infty$ to $\infty$, each denominator factor has argument
increasing from $0$ to $\pi$. The resulting line direction traverses
all angles modulo $\pi$. Choose it parallel to the real line for a
parabola, or at the angle between the two asymptotes for a hyperbola.
If a boundary point equals the selected pole, perturb that boundary
point on the real line, keeping $u,v,t$ fixed. Other poles in the
partner are handled by the same approximation. Thus it suffices to
treat five points distributed on two lines with the prescribed angle.

For a hyperbola move the lines by a similarity to its two asymptotes.
In real linear coordinates $(x,y)$ along the asymptotes the fixed
hyperbola has equation $xy=k\ne0$. Its uniform dilates $xy=\epsilon^2k$
approximate any finite configurations on the union of the axes:
for example $(x,\epsilon^2 k/x)\to(x,0)$ at $x\ne0$, with the
other axis handled by exchanging coordinates. A point at the origin
is also a limit, for example $(\epsilon,\epsilon k)$. Use the same
dilation for all five points. Pull each approximating pair back to the
fixed hyperbola by the inverse similarity before using the hypothesis.

For a parabola all nondegenerate real parabolas are similar to
$p(t)=t+it^2$. The upper line has positive height, so a similarity
takes the two parallel lines to $\Re z=\pm1$. For a target
$\eta+iy$, $\eta=\pm1$, choose $s=y/(2\eta)$. Then
\[
 [p(\eta T+s)-iT^2]/T=\eta+iy+(s+i s^2)/T.
\]
These points approach the entire specified configuration using a
common similarity. In both cases take the limit at a fixed normalized
partner, undo the chart, and use closedness to include all degenerate
input lists.
\end{proof}

In case F of the question the literal exclusion of the spin image does
not restrict the conjugated second matrix: multiplication of $W$ by
a unit scalar can take its determinant away from one. If the exclusion
instead concerns the entire rephasing and permutation orbit of the spin
image, that orbit already has the separate spin certificate in the
research report (Section 8). The general inclusion on these conics
remains unresolved.
\end{document}
